Put $A=(0,0),C=(7,0)$. The two side lengths give $B=(\dfrac{33}{14},\dfrac{15\sqrt3}{14})$. Each circle centre lies vertically over its tangency point. From $AB^2=2O_1A\cdot y_B$ and $CB^2=2O_2C\cdot y_B$, their heights are $\dfrac{21}{(5\sqrt3)}$ and $\dfrac{35}{(3\sqrt3)}$. The quadrilateral is a trapezoid of width seven, so its area is $(\dfrac{7}{2})(\dfrac{21}{(5\sqrt3)}+\dfrac{35}{(3\sqrt3)})=\dfrac{833}{(15\sqrt3)}$. Thus $\boxed{a+b+c=851}$.
\begin{asy}
import graph;import geometry;usepackage("amsmath");size(7cm);
pair A=(0,0),C=(7,0),B=(33/14,15sqrt(3)/14),O1=(0,21/(5sqrt(3))),O2=(7,35/(3sqrt(3)));draw(A--B--C--cycle);draw(A--O1--O2--C,blue);draw(circle(O1,O1.y),gray);draw(circle(O2,O2.y),gray);label("$A$",A,(-1,-1));label("$B$",B,(0,1));label("$C$",C,(1,-1));label("$O_1$",O1,(-1,0));label("$O_2$",O2,(1,0));
\end{asy}

